\section{Limitations and Future Work}\label{sec:limitations}

The three correctness properties established in Section~\ref{sec:correctness}
are conditional guarantees: each holds under a defined set of structural and
deployment assumptions.  This section makes those assumptions explicit,
characterizes the operational consequences of their failure, and maps each
gap to a directed research agenda.

% --------------------------------------------------------------
\subsection{Limitations}\label{sec:limitations:subsection}
% --------------------------------------------------------------

\subsubsection{Human Response Latency}\label{sec:human-response-latency}

The \bxthree{} Framework's Liveness property guarantees eventual human contact
within a bounded wall-clock interval, but the bound $T_{\max}$ is a function of
network topology, retry parameters, and provisioned timeouts --- not of human
cognitive response time.  For operational domains with time constants measured
in milliseconds (high-frequency trading, real-time process control, autonomous
vehicle navigation), the time required for a human principal to receive,
interpret, and respond to an escalation may exceed the system's operational
time horizon by one to two orders of magnitude.

The framework's prescribed failure mode under latency stress is well-defined:
the node stalls in \texttt{ESCALATING} until the anchor responds or all
pathways time out.  This is the intended safety-preserving behavior.  However,
the operational cost of a stall is domain-dependent in ways the framework
does not currently reason about.  In a precision agriculture context (Irrig8
deployment), a stall that delays a watering decision by minutes is acceptable;
in a financial risk-interlock context, a stall that delays a margin-call
response by even seconds may expose the system to catastrophic loss.  The
framework treats the human anchor as a timing black box; there is no mechanism
to reason about whether delay reflects anchor overload, a structural mismatch
between $T_{\max}$ and the anchor's operational availability, or a pathway
failure.

The deeper limitation is that human response latency is provisioned statically
at node setup and is not sensitive to the anchor's current operational context.
$T_{\max}$ is a deployment parameter, not an adaptive runtime parameter.  For
latency-sensitive deployments, the provisioning burden falls on the system
designer to establish empirically grounded latency budgets and encode them as
enforceable policy.

\textit{Directed research:} Empirical characterization of human response
latency distributions across deployment contexts to establish statistically
grounded, domain-specific latency budgets.  Integration of these budgets into
the node-provisioning interface as enforceable policy parameters with runtime
monitoring.  Investigation of graded escalation responses: a protocol that
requests a rapid \texttt{ACK} followed by a detailed directive would better
match the operational reality of human decision-making under varying loads.

\subsubsection{Adversarial Conditions}\label{sec:adversarial-conditions}

The Safety property requires that human-contact pathways cannot be permanently
closed without recorded human authorization.  This guarantee holds under the
assumption that the host platform provides independent scheduling priority to
the Bailout Monitor (Assumption A1.1).  In an adversarial environment ---
particularly one in which the host platform itself may be compromised --- this
assumption may not hold.  A sophisticated attacker controlling the execution
substrate can modify thread priorities, preempt the monitor thread, introduce
race conditions into pathway-closure operations, or rewrite the Bailout
Monitor's code before it executes.

The framework's guarantees are platform guarantees.  They rest on invariants
that the host platform must supply: non-preemptible monitor scheduling,
append-only persistence for the authorization ledger, and interruptible
blocking calls for pathway management.  The \bxthree{} Framework is designed
for production AI deployments in which individual agents may behave
unexpectedly but the host platform is trusted --- the standard assumption for
regulated industries.  It is not designed for deployments in which the platform
itself is an active adversary: national infrastructure control systems,
military command-and-control environments, financial clearinghouse
infrastructure.

\textit{Directed research:} Formal specification of the composition interface
between the \bxthree{} Framework and substrate-integrity mechanisms.
Investigation of whether the framework's guarantees can be preserved under
weaker platform assumptions, including partially compromised substrates.
Exploration of hardware-attested execution environments (e.g., Intel SGX,
ARM TrustZone) as a deployment target for the Bailout Monitor, eliminating
the OS as a trust dependency.

\subsubsection{Scalability}\label{sec:scalability}

The \bxthree{} Framework is specified with a 1:1 relationship between each node
and its human anchor.  This model is appropriate for single-node deployments or
small-scale systems, but it creates a structural bottleneck in enterprise
contexts where one anchor may supervise hundreds of concurrent nodes.  Without
pooling and load-balancing mechanisms, the human anchor becomes the
rate-limiting step in a system whose design intent is to route exceptions away
from centralized machine decision points.

The current specification treats human anchoring as a static 1:1 relationship
established at node provisioning.  This model must evolve to support at minimum
1:N relationships (one anchor supervising multiple nodes with workload pooling)
and, in advanced deployments, N:M relationships (anchor pools with load
balancing and dynamic routing).  The correctness properties of the framework
must be shown to extend to these pooled configurations, particularly Safety:
no pooling configuration may produce a state in which all human-contact
pathways for a given node are permanently disabled.

Additional scalability considerations arise from the spawn tree depth.  As the
depth of the accountability chain grows through recursive spawning in large
multi-agent systems, the time to propagate an exception to the anchor grows
proportionally.  The interaction between chain depth, propagation latency,
and timeout configuration must be characterized for deployments with large
spawn trees.

\textit{Directed research:} Formal specification of human anchor pool
architectures, including load-balancing policies, prioritization schemes, and
routing algorithms.  Extension of Safety and Liveness proofs to pooled anchor
configurations with formally verified pathway-redundancy guarantees.
Empirical evaluation of anchor pool performance under simulated
high-escalation-load conditions.

% --------------------------------------------------------------
\subsection{Future Work}\label{sec:future-work}
% --------------------------------------------------------------

The following research directions extend the \bxthree{} Framework's guarantees
to broader deployment contexts and establish a principled research program
beyond the current specification.

\textbf{Formal verification.}  Full mechanical verification of the Safety,
Liveness, and Accountability invariants using a formal verification tool
(Coq, Isabelle/HOL, or TLA\textsubscript{+}) applied to the \bxthree{} three-layer
model and all five pillar state machines.  The current paper establishes the
invariants through proof sketches; the directed research agenda targets full
formal verification as a precondition for deployment in regulated industries.

\textbf{Adaptive timeout mechanisms.}  Replacement of static timeout
parameters ($T_{\text{bounds}}$, $T_{\text{prop}}$, $T_{\text{human}}$) with
adaptive mechanisms that adjust to anchor workload, pathway latency, and
deployment context.  The adaptive protocol would maintain the safety invariant
(the anchor is always contacted within $T_{\max}$) while reducing unnecessary
stalls in low-consequence, high-latency scenarios.

\textbf{Distributed human anchor pools.}  Formal specification and empirical
evaluation of anchor pool architectures that sustain scalability under
cascading failure conditions without compromising the no-machine-actor-absorbs
property.  This is the primary scalability extension required for enterprise
\bxthree{} deployments.

\textbf{Exploitation resistance.}  Design of triggering-behavior anomaly
detectors that audit the statistical relationship between observed exception
conditions and triggered escalations per node over rolling observation windows.
Investigation of provenance attestation for operating-range assessments:
requiring that each Bounds Engine self-assessment be counter-signed by an
independent audit process, making post-hoc misclassification detectable through
forensic reconstruction.